%
    \subsection{Propagator-inspired markers of epileptogenic tissue}
        \label{ssec:epileptogenic}
%
\paragraph{Seizure-onset contacts form a co-organized diffusion group, not a contact-by-contact label.}
The cognitive trace is a cross-phase correlation --- it measures how much the task's reorganization of the diffusion hierarchy persists into rest. That hierarchy is built from a propagator on the connectivity graph, the heat kernel \(e^{-\tau L}\); read directly instead, within a single recording, as the diffusion affinity between two contacts, the same propagator localizes epileptogenic tissue. A contact-by-contact marker of the seizure-onset zone does not exist here: within a patient any contact-level signal collapses into how strongly the contact is coupled, and across patients no contact-level rule transfers. Read relationally, the picture changes. The seizure-onset contacts diffuse together: a contact's heat-kernel affinity to the rest of the seizure set separates seizure from healthy contacts, over and above node strength --- the summed weight of a contact's connections, which by itself carries some signal because seizure tissue tends to be hyperconnected. With that strength baseline removed, the affinity still reaches \gls{auc} \(0.83\) in \(\delta\), \(0.82\) in \(\lowgamma\), and \(0.745\) in \(\beta\) --- higher, in each band, than seizure sets of matched coupling strength reach by chance --- clearing that matched-strength null in \(7\) of \(10\) patients for \(\delta\), \(8\) for \(\lowgamma\), and \(8\) for \(\beta\), strongest in \(\delta\) and present across several rhythms, with \(\alpha\) marginal (\(0.61\), \(5\) of \(10\); \FigRef{fig:epi1}b). Reading the affinity across diffusion scales rather than at the single fine scale sharpens this separation in most patients (multiscale above single-scale in \(7/10\) for \(\delta\), \(8/10\) for \(\beta\), \(7/10\) for \(\lowgamma\)), so the diffusion depth \(\tau\) improves how well seizure tissue ranks. These values reproduce, and slightly exceed, the fully connected marker once the connectivity graph is sparsified to its backbone --- the localization is a property of the diffusion operator itself, not of leaving the graph dense. Because the \gls{imcoh} substrate is immune to zero-lag volume conduction, this grouping cannot be spatial proximity --- and given seizure contacts on some electrodes as seeds, the same affinity ranks the seizure contacts on \emph{other} electrodes above healthy ones (\(\delta\), \gls{auc} \(0.72\), \(8\) of \(10\) patients, a shuffle of the labels collapsing to chance), the non-trivial case in which distance to a known contact is useless (\FigRef{fig:epi1}a). The seizure zone is thus organized as a strength-independent, distributed diffusion group, not a set of individually marked contacts.

\paragraph{A calibrated seizure-zone probability, carried by \(\delta\) and sharpened across rhythms.}
Turned into a deployable marker --- trained on some patients and applied to the next --- the single strongest band is \(\delta\): from three known seizure contacts as seeds, a logistic model on its diffusion features alone separates seizure from healthy contacts at a mean \gls{auc} of \(0.76\) across patients. Fusing all six bands raises that separation only slightly, to \(0.81\), but it nearly doubles the precision at the top of the ranking, where a shortlist is drawn: the five highest-scoring contacts per patient go from \(34\%\) to \(60\%\) seizure-onset, roughly a sevenfold enrichment over the base rate. The discrimination is therefore carried by \(\delta\); the other rhythms add largely independent information that sharpens the top of the list rather than the overall separation. The diffusion-scale sweep and the band fusion act on different axes: sweeping \(\tau\) lifts each band's ranking (\gls{auc}) but not its shortlist --- single-band precision@\(5\) stays near \(0.40\) at every scale --- so the near-doubling of top-five precision comes from fusing the bands, not from the scale sweep. The fused model is the one we report: leave-one-patient-out, it assigns every contact a calibrated probability of seizure onset --- median \gls{auc} \(0.87\), probability \(0.38\) at true seizure contacts against \(0.08\) at healthy ones, \(9\) of \(10\) patients above chance, and a label-shuffle null at \(0.48\) (\FigRef{fig:epi1}c). A nonlinear model reaches a higher \gls{auc} (\(0.86\)) with no gain in precision and a clear overfitting risk on ten patients, so we keep the interpretable logistic. No contact coordinates enter it --- only diffusion features on the volume-conduction-immune graph --- so the probability is a connectivity read-out, not a restatement of where the electrodes were placed.

%
\begin{figure*}[p]
    \centering
    \includegraphics[width=\textwidth,height=0.95\textheight,keepaspectratio]%
        {fig_epi1.pdf}

    \caption{\textbf{The same diffusion propagator that carries the cognitive trace
        localizes epileptogenic tissue.}
        \textbf{(a)} The seizure zone recovered as a spreading process (exemplar
        \patient{08}, \(\delta\), post-task rest). Heat placed on the known
        seizure-onset contacts of one electrode (gold diamonds) diffuses across the
        volume-conduction-immune \gls{imcoh} graph through the heat kernel
        \(e^{-\tau L}\); every other contact is colored and sized by the resulting
        seed affinity (strength-residualized rank) and the warmest flows are drawn as
        arcs. The heat reaches the \emph{unseeded} seizure-onset contacts on a
        \emph{different} electrode (red rings) --- distant discovery with proximity
        removed by construction --- while healthy cortex stays cool.
        \textbf{(b)} Read relationally over all contacts, the strength-residual
        affinity separates seizure from healthy contacts above node strength
        (\(0.5\) = strength alone): cohort-median multiscale \gls{auc} \(0.83\) in
        \(\delta\), \(0.82\) in \(\lowgamma\), \(0.745\) in \(\beta\), \(0.61\) in
        \(\alpha\), each above a matched-strength fake-seizure null, clearing it in
        \(7\)/\(8\)/\(8\)/\(5\) of \(10\) patients (green = beats own null;
        \(0.5\) = strength alone / chance). Reading across diffusion scales (filled)
        rather than at the single fine scale (open) raises the ranking in the
        majority of patients (\(\delta\) \(7/10\), \(\beta\) \(8/10\), \(\lowgamma\)
        \(7/10\)).
        \textbf{(c)} Fused across the six bands into a leave-one-patient-out logistic
        detector, the propagator features yield a calibrated seizure-onset
        probability. Fusion lifts precision@5 from \(34\%\) (\(\delta\) alone) to
        \(60\%\) (\(\approx\!7\times\) the base rate) while adding little \gls{auc} (a
        nonlinear model buys \gls{auc} but no precision; grey marker); the reliability
        curve (right) tracks the diagonal, with median \gls{auc} \(0.87\), \(9/10\)
        patients above chance, and true-onset probability \(0.38\) against \(0.08\) at
        healthy contacts.}
    \label{fig:epi1}
\end{figure*}
%
\paragraph{The cohort splits into two seizure-network shapes.}
The cohort mean hides two populations. In eight patients the seizure zone is a co-diffusing community and the detector is strong --- \gls{auc} up to \(0.99\) (five of the eight above \(0.90\)), with all five top-ranked contacts seizure-onset in the clearest cases. In the other two, both right-hemisphere implants, the seizure zone \emph{is} a network hub rather than a community, and the propagator read-out fails (\gls{auc} \(0.49\) and \(0.57\)); there, plain coupling strength predicts the seizure contacts the propagator misses (in one, strength reaches \gls{auc} \(0.93\) where the propagator sits near chance). The split is the same across every read-out --- the relational grouping, the distant-seed discovery, and the fused detector all succeed on the community patients and fail on the two hub patients --- so it is a property of the seizure network's shape, not of the method. Choosing per patient between the affinity and the strength read-out for the hub-type implants recovers nine of ten; in the tenth, both read-outs fail (\FigRef{fig:epi2}a). Performance is therefore reported per patient rather than as a pooled mean, which would average a working detector together with a known failure mode.

\paragraph{A mechanistic marker of the seizure-onset zone, calibrated for triage.}
The marker is mechanistic and coordinate-free: it is computed only from the diffusion propagator on the volume-conduction-immune connectivity graph, carries no amplitude, tissue, or proximity prior, and transfers across patients --- a connectivity signal that co-locates with the clinically defined seizure-onset zone, read from the same operator that carries the cognitive trace. Because it is calibrated against those clinical labels rather than against resection and surgical outcome, it is a triage aid --- a top-five precision of \(60\%\), a roughly sevenfold enrichment over the base rate --- rather than a standalone localizer. It applies to gray and white matter alike: \(41\%\) of the seizure-onset contacts lie in white matter, which the marker ranks without a gray-matter prior (\FigRef{fig:epi2}b). Its highest-ranked unmarked contacts --- in hippocampus, amygdala, and the medial orbitofrontal cortex where the \(\beta\) trace concentrates --- are candidate occult sites for prospective testing, any of which, if genuine, would raise the precision reported here.

%
\begin{figure*}[p]
    \centering
    \includegraphics[width=\textwidth,height=0.86\textheight,keepaspectratio]%
        {fig_epi2.pdf}

    \caption{\textbf{Two seizure-network shapes, and the honest ceiling of the marker.}
        \textbf{(a)} Per-patient detector \gls{auc} (left) splits the cohort in two:
        eight patients whose seizure zone is a co-diffusing community (green,
        \gls{auc} up to \(0.99\), five above \(0.90\)) and two right-hemisphere
        implants whose seizure zone is a network hub (red, \gls{auc} \(0.49\) and
        \(0.57\)), where the propagator fails and node strength predicts the onset
        contacts instead. Selecting per patient between the affinity and the strength
        read-out (right; open \(\to\) filled marker) rescues the hub implants and
        recovers \(9/10\) (cohort off-shaft \gls{auc} \(0.72 \to 0.75\)); \patient{10}
        is unrecoverable, both read-outs failing.
        \textbf{(b)} The marker is a triage aid, not a standalone localizer.
        Precision@\(k\) (left) is \(60\%\) for the top-five shortlist
        (\(\approx\!7\times\) the per-patient base-rate band, grey) and relaxes toward
        the base rate as the list lengthens --- a ceiling set by target rarity, not by
        the ranker. It spans tissue classes: \(41\%\) of seizure-onset contacts lie in
        white matter, which the marker ranks without a gray-matter prior. Its
        highest-ranked \emph{unmarked} contacts (right) --- in hippocampus, amygdala,
        and the medial orbitofrontal cortex where the \(\beta\) trace concentrates
        (blue, orange) --- are candidate occult sites for prospective testing
        (hypotheses, none outcome-validated).}
    \label{fig:epi2}
\end{figure*}